\DocSection{Exact Validation Framework}
\label{sec:validation}

Correct implementation of an MCMC transition does not establish that the
intended density was specified.  Validation therefore starts with finite MDPs
for which all \(|\mathcal A|^D\) stationary deterministic policies can be
enumerated, where \(D\) is the number of decision states.  Finite-horizon
dynamic programming computes \(\ExpectedReturn_H(\policy)\) for every policy,
and log-sum-exp normalization gives the exact target probabilities, partition
function, and action marginals in Equation~\eqref{eq:policy-target}.

\DocSubsection{Target-substitution diagnostic}

The smallest diagnostic contains two policies with equal expected return but
different return variance.  Under a uniform reference measure,
Equation~\eqref{eq:policy-target} assigns them equal probability.  A method
that averages \(\exp(\beta\Return_H)\) instead assigns more mass to the
variable policy.  This is a deterministic pass-or-fail test of the order of
operations rather than a performance benchmark.

Additional exact tests cover state revisitation, absorbing states, finite
horizons, policy memoization, common random numbers, Metropolis--Hastings
proposal ratios, replica swaps, particle ancestry, and score-function
gradients.  In particular, empirical frequencies of both single-chain and
replica-exchange samplers must agree with the enumerated target on tiny MDPs.

\DocSubsection{Fail-fast gates}

Experiments advance only after the preceding gate passes:
\begin{enumerate}[label=G\arabic*.]
  \item \textbf{Mechanics.} Transition rows, rewards, termination, horizons,
        seeded replay, artifact writing, and budget counters pass unit tests.
  \item \textbf{Target.} Enumeration verifies exact probabilities, modes,
        marginals, local acceptance ratios, and swap acceptance ratios.
  \item \textbf{Approximation.} Increasing the immutable tape-bank size
        reduces sample-average target error on repeated small problems; a
        failure stops simulator-only scaling.
  \item \textbf{Mixing.} One-seed medium pilots show finite outputs, adequate
        cold-chain effective sample size, and exchange across every ladder
        edge.  A high local acceptance rate alone is not sufficient.
  \item \textbf{Cost.} Three short runs establish memory, wall-clock time, and
        value-evaluation or simulator-step rates before a large job array is
        submitted.
\end{enumerate}

The final thresholds for effective sample size and target error are frozen
after pilot calibration and before confirmatory seeds are inspected.  Failed,
timed-out, and numerically invalid runs remain in the accounting.
